
Eine reale Spannungsquelle kann durch eine ideale Spannungsquelle mit der Leerlaufspannung $U_0$ und einem in Serie 
geschalteten Innenwiderstand $R_I$ modelliert werden. An diese Quelle ist ein variabler Lastwiderstand $R_L$ angeschlossen.
Mit dieser Schaltung kann anschliessend die maximal abgegebene Leistung berechnet werden.
\begin{figure}[H]
    \centering
    \includegraphics[width=1\textwidth]{Spannungsquelle.png}
    \caption{Leistungsanpassung mit Spannungsquelle}
    \label{fig:Spannungsquelle}
\end{figure}
\noindent
Die Werte der Oben dargestellten Schaltung betragen:
\begin{equation*}
    U_0 = 15V
\end{equation*}
\begin{equation*}
    R_i = 5\Omega
\end{equation*}